\documentclass[10pt]{article}
\usepackage[T1]{fontenc}
\usepackage[utf8]{inputenc}
\usepackage[french]{babel}
\usepackage{amsmath,amssymb,mathrsfs,mathpazo}
\usepackage{geometry}
\geometry{a4paper,margin=2cm}
\usepackage{graphicx,array,tabularx,tabularray,colortbl,multirow,multicol,booktabs,enumitem,pifont,ulem,fancybox,diagbox,cancel}
\usepackage{pstricks,pstricks-add,pst-plot,pst-tree,pst-node,pst-text,pst-func,pst-math,pst-bspline,pst-3dplot,pst-eucl,pst-solides3d}
\let\pscalculate\undefined
\usepackage{pst-bezier}
\usepackage{tikz,pgfplots,tkz-tab}
\usetikzlibrary{arrows,arrows.meta,patterns,calc,positioning,trees,decorations.markings,shapes.geometric,shapes.misc,shapes.arrows,intersections,decorations.pathmorphing,decorations.pathreplacing,backgrounds,shadows,babel}
\pgfplotsset{compat=1.18}
\usepgfplotslibrary{fillbetween}
\usepackage[np]{numprint}
\usepackage{listings,makecell,pas-tableur,twoopt}
\let\barre\undefined
\usepackage[pstricks]{tablvar}
\definecolor{darkgreen}{RGB}{0,100,0}
\providecommand{\red}{\color{red}}
\providecommand{\blue}{\color{blue}}
\providecommand{\green}{\color{green}}
\providecommand{\black}{\color{black}}
\providecommand{\white}{\color{white}}
\providecommand{\gray}{\color{gray}}
\providecommand{\cyan}{\color{cyan}}
\providecommand{\magenta}{\color{magenta}}
\providecommand{\yellow}{\color{yellow}}
\providecommand{\violet}{\color{violet}}
\providecommand{\orange}{\color{orange}}
\providecommand{\euro}{\texteuro}
\providecommand{\EUR}{\texteuro}
\providecommand{\ieme}{\textsuperscript{e}}
\providecommand{\degres}{\ensuremath{^\circ}}
\providecommand{\degree}{\ensuremath{^\circ}}
\providecommand{\vv}[1]{\overrightarrow{#1}}
\providecommand{\Coord}[1]{#1}
\providecommand{\decosix}{}
\providecommand{\starredbullet}{\ensuremath{\bullet}}
\pagestyle{empty}
\setlength{\parindent}{0pt}
\usepackage{adjustbox}\DeclareUnicodeCharacter{2236}{:}\DeclareUnicodeCharacter{2190}{\ensuremath{\leftarrow}}
\providecommand{\R}{\mathbb{R}}
\providecommand{\N}{\mathbb{N}}
\providecommand{\Z}{\mathbb{Z}}
\providecommand{\Q}{\mathbb{Q}}
\providecommand{\D}{\mathbb{D}}
\providecommand{\C}{\mathbb{C}}
\providecommand{\e}{{\mathrm{e}}}
\providecommand{\vect}[1]{\overrightarrow{#1}}
\providecommand{\vectt}[1]{\overrightarrow{\mathrm{#1}}}
\providecommand{\vvt}[1]{\overrightarrow{\mathrm{#1}}}
\providecommand{\np}[2][]{{#2~#1}}
\providecommand{\vv}[1]{{\overrightarrow{#1}}}
\renewcommand{\d}{{\mathrm{d}}}
\providecommand{\degres}{{^\circ}}
\providecommand{\degre}{{^\circ}}
\providecommand{\blue}{\color{blue}}
\providecommand{\red}{\color{red}}
\providecommand{\enskip}{\quad}
\providecommand{\llbracket}{\lbrack\!\lbrack}
\providecommand{\rrbracket}{\rbrack\!\rbrack}
\providecommand{\euro}{{\text{€}}}
\providecommand{\no}{{\text{n°}}}
\providecommand{\interval}[2]{{[#1;#2]}}
\providecommand{\intervFO}[2]{{[#1;#2[}}
\providecommand{\dotfill}{\ldots}

\providecommand{\ds}{\displaystyle}
\providecommand{\textding}[1]{#1}
\providecommand{\Oij}{\(\mathrm{O};\vec{\imath},\vec{\jmath}\)}
\providecommand{\Oijk}{\(\mathrm{O};\vec{\imath},\vec{\jmath},\vec{k}\)}
\providecommand{\Ouv}{\(\mathrm{O};\vec{u},\vec{v}\)}
\providecommand{\intervFF}[2]{\ensuremath{[#1;#2]}}
\renewcommand{\euro}{\ifmmode\text{\texteuro}\else\texteuro\fi}
\DeclareUnicodeCharacter{00B0}{\ifmmode^{\circ}\else\textdegree\fi}
\DeclareUnicodeCharacter{20AC}{\ifmmode\text{\texteuro}\else\texteuro\fi}
% Énoncés : Éducation nationale. Transcription et corrigés : APMEP, Denis Vergès et contributeurs. Découpage et classement : CoopMaths, Sébastien Lozano, Éric Elter et contributeurs. Adaptation de la présentation : Maths Explorer.
% CC BY-SA : https://coopmaths.fr/www/about/
% https://coopmaths.fr/alea/static/bac/2021/tex/bac_2021_06_asie_7.tex
\begin{document}
\noindent\textit{APMEP / CoopMaths. CC BY-SA. Présentation adaptée par Maths Explorer.}\par\medskip
\textbf{Commun à tous les candidats}

\medskip

On considère un pavé droit ABCDEFGH tel que AB = AD = 1 et AE = 2, représenté ci- dessous.

Le point I est le milieu du segment [AE]. Le point K est le milieu du segment [DC]. Le point L
est défini par: $\vect{\text{DL}} = \dfrac{3}{2}\vect{\text{AI}}$. N est le projeté orthogonal du point D sur le plan (AKL).

\begin{center}
\psset{unit=1cm,arrowsize=2pt 4}
\begin{pspicture}(6,10)
\pspolygon(0.5,0.5)(4,0)(4,7.2)(0.5,7.7)
\uput[d](0.5,0.5){B} \uput[d](4,0){C} \uput[ul](4,7.2){G} \uput[ul](0.5,7.7){F} 
\psline(4,0)(5,1.4)(5,8.6)(4,7.2)
\uput[r](5,1.4){D} \uput[ur](5,8.6){H} \uput[d](1.5,1.9){A} \uput[dr](4.5,0.7){K}\uput[l](1.5,5.5){I}
\uput[r](5,7.3){L}
\psline(5,8.6)(1.5,9.1)(0.5,7.7)
\uput[u](1.5,9.1){E}
\psline[linestyle=dashed]{->}(1.5,1.9)(1.5,5.5)
\psline[linestyle=dashed]{->}(1.5,1.9)(0.5,0.5)
\psline[linestyle=dashed]{->}(1.5,1.9)(5,1.4)
\psline[linestyle=dashed](4.5,0.7)(1.5,1.9)(5,7.3)
\psline[linestyle=dashed](1.5,5.5)(1.5,9.1)
\psline(4.5,0.7)(5,7.3)
\end{pspicture}
\end{center}

\bigskip

On se place dans le repère orthonormé $\left(\text{A}~;~\vect{\text{AB}},~\vect{\text{AD}},~\vect{\text{AI}}\right)$. 

On admet que le point L a pour coordonnées $\left(0~;~1~;~\dfrac{3}{2}\right)$.

\medskip

\begin{enumerate}
\item Déterminer les coordonnées des vecteurs $\vect{\text{AK}}$ et $\vect{\text{AL}}$.
\item  
	\begin{enumerate}
		\item Démontrer que le vecteur $\vect{n}$ de coordonnées $(6~;~-3~;~2)$ est un vecteur normal au plan (AKL).
		\item En déduire une équation cartésienne du plan (AKL).
		\item Déterminer un système d'équations paramétriques de la droite $\Delta$ passant par D et
perpendiculaire au plan (AKL).
		\item En déduire que le point N de coordonnées $\left(\dfrac{18}{49}~;~\dfrac{40}{49}~;~\dfrac{6}{49}\right)$ est le projeté orthogonal du point D sur le plan (AKL).
	\end{enumerate}
\end{enumerate}

On rappelle que le volume $\mathcal{V}$ d'un tétraèdre est donné par la formule : 

\[\mathcal{V} = \dfrac{1}{3}\times  (\text{aire de la base}) \times \text{hauteur}.\]

\begin{enumerate}[resume]
\item 
	\begin{enumerate}
		\item Calculer le volume du tétraèdre ADKL en utilisant le triangle ADK comme base. 
		\item Calculer la distance du point D au plan (AKL).
		\item Déduire des questions précédentes l'aire du triangle AKL.
	\end{enumerate}
\end{enumerate}

\bigskip

\end{document}
